\section{System Architecture}
\label{sec:architecture}

\subsection{Repository layout}
\label{sec:layout}

The repository (Figure~\ref{fig:layout}) contains two separately installable packages and a
published example run. The \texttt{main} branch is the clean, publishable surface. A \texttt{dev}
branch preserves the full research history: the research log from which \repo{FINDINGS.md} is
curated, standalone study drivers and experiment configurations that reproduce individual findings.

\begin{figure}[t]
\centering
\footnotesize
\begin{minipage}{0.94\textwidth}
\renewcommand*\DTstyle{\footnotesize\ttfamily}
\renewcommand*\DTstylecomment{\footnotesize\rmfamily\itshape\color{gray!60!black}}
\setlength{\DTbaselineskip}{11pt}
\dirtree{%
.1 TSFM-Interp/.
.2 README.md, FINDINGS.md, DEPENDENCIES.md, CONTRIBUTING.md, CITATION.cff.
.2 tsfm\_benchmark/\DTcomment{corpus generation and validation (installable package)}.
.3 build\_pipeline/\DTcomment{generators, leakage audit, sealing, corruptions, counterfactuals}.
.3 benchmark\_validation/\DTcomment{catch22 diversity, redundancy matching, split exchangeability}.
.3 configs/, example\_runs/, tests/\DTcomment{golden-hash regeneration tests}.
.2 tsfm\_lens/\DTcomment{model analysis (installable package)}.
.3 run.py, configs/examples/, docs/, tests/\DTcomment{CLI, presets, adapters guide, 171 test modules}.
.3 tsfm\_lens/.
.4 models/\DTcomment{adapters, derived tiers, conformance, generic\_hf, template}.
.4 extraction/\DTcomment{hooks, span discovery, alignment gate, zarr store}.
.4 analysis/\DTcomment{L0--L3, lens, attention, statistics, register/confirm}.
.4 sae/\DTcomment{TopK SAE, ablation battery, atlas, transfer, shared-input agreement}.
.4 report/\DTcomment{section builders, derived scorecard, findings, sanitizer}.
.4 pipeline.py, manifest.py, config.py\DTcomment{stages, gates, fingerprints}.
.2 examples/concept\_atlas\_v2/\DTcomment{report.html, config.yaml, run/ (per-stage artifacts)}.
.2 paper/\DTcomment{this paper; figures/make\_figures.py regenerates every data figure}.
}
\end{minipage}
\caption{Repository layout of the \texttt{main} branch. The benchmark and analysis packages are
independent; the example run folder contains every lightweight artifact of the run analyzed in
Section~\ref{sec:casestudy}, so its report and this paper's figures can be regenerated without GPUs.}
\label{fig:layout}
\end{figure}

\subsection{Controlled benchmark generation}
\label{sec:benchmark}

\paragraph{Two leakage tiers.} Real data enter the benchmark only as leakage references, as realism
calibration and as raw material for explicitly labeled generators. The \emph{synthetic} tier touches
no real data and has exact ground truth: the \texttt{parametric} generator is additive (trend,
multiple seasonalities, AR noise, changepoints, anomalies, random walk, intermittency and
heteroskedasticity), and \texttt{random\_parametric} draws a recipe from one of eight frozen
archetypes (trend-dominant, seasonal-dominant, multi-seasonal, regime-switching, AR colored noise,
anomaly-heavy, clean low-noise, noisy chaotic), with four more available opt-in. The
\emph{real-derived} tier (\texttt{mixture}, \texttt{block\_bootstrap} \citep{kunsch1989block} and
\texttt{sequential\_par}) inherits its source distribution, and its ground truth is attributional (which
source, which blocks). Every sample is labeled with its tier. Mixing the tiers in one corpus is a
deliberate design decision: ``no leakage'' and ``realistic'' pull in opposite directions, and the
labels let each analysis report its results per tier rather than choosing one of the two.

\paragraph{Leakage audit and sealing.} Every candidate sample passes a two-stage leakage gate against
real reference collections (the Monash archive \citep{godahewa2021monash}): a cheap prefilter,
then banded dynamic time warping \citep{sakoe1978dtw}. It must also pass a cross-split near-duplicate
check. Dev and private splits draw from disjoint seed ranges; seeds are mixed with SHA-256 (never a
process-salted hash), and every sample is hashed so that \texttt{load\_sealed(verify=True)} detects
tampering. Fresh private \emph{epochs} are minted by striding the seed space, so a consumed private
split can be replaced without touching the dev split.

\paragraph{Validation.} A separate package checks that a corpus is diverse and non-redundant.
Redundancy is measured by shape matching (cross-correlation or DTW). Diversity is measured in the
catch22 feature space \citep{lubba2019catch22}, never on a 2-D embedding, using the participation-ratio
effective dimension \citep{gao2017pr}, nearest-neighbor tails and near-collision rates overall and per
group. Split exchangeability, on which the confirmatory design depends, is tested with an energy
distance \citep{szekely2013energy} plus per-feature TOST equivalence \citep{schuirmann1987tost}, with a
three-state verdict. On the v1 benchmark, dev and private are exchangeable (energy $p=0.51$; 21 of 22
features equivalent; \fid{MP-09}).

\paragraph{Controlled perturbations.} Two perturbation families are built in. Nine structured
\emph{corruptions} (noise, detrend, deseasonalize, frequency shift, level shift, spike, smooth, warp,
dropout) drive the L3 sensitivity and patching analyses; because the underlying components are known,
their effect on the input is known exactly. \emph{Generator counterfactuals} regenerate a series with
one generative knob moved along a dose ladder while every random draw stays fixed; dose~1.0 must
reproduce the original bit-exactly, a free identity test that caught a real bug (\fid{BM-05}).

\subsection{Model adapters and capability tiers}
\label{sec:adapters}

A model enters the system through an adapter (\repo{tsfm_lens/tsfm_lens/models/base.py}) that exposes
\texttt{load}, \texttt{predict} (point forecast plus requested quantiles, in the original scale), the
underlying module, input preparation, a forward pass and \texttt{token\_time\_spans} (one contiguous
time interval per token). Optional methods expose attention and MLP modules and attention patterns.
The adapter's \emph{capability tier} is derived from which methods a subclass actually overrides, never
declared: tier~0 (black box: load and predict), tier~1 (observable: hidden states with spans), tier~2
(steerable: single-pass context, so patching is possible) and tier~3 (decomposable: attention internals).
A stage that needs a higher tier than a model has is skipped for that model, and the report names the
missing method. A \emph{missing capability} raises \texttt{CapabilityUnavailable}, a fact about the
adapter; a \emph{measured} failure of time localization raises \texttt{NotTimeLocalized}, a fact about the
model, and routes it to behavioral analysis only.

Adding a model takes one of three routes. The zero-code \texttt{generic\_hf} adapter discovers layers
and measures spans for a Hugging Face checkpoint and refuses rather than guesses. It has been verified
on Timer, which it recovers with 96-step spans from measurement alone (\fid{MP-01}). A hand-written
adapter is a copy of \repo{tsfm_lens/tsfm_lens/models/TEMPLATE_adapter.py} placed in
\texttt{models/contrib/}, where it is discovered by an AST scan. In either case
\texttt{run.py -{}-check-adapter} runs a one-command checklist: patch-identity and patch-reach
conformance, the alignment gate, span discovery and capability rows. The adapters guide
(\repo{tsfm_lens/docs/ADDING_A_MODEL.md}) documents the pitfalls we encountered, including a worked
account of an adapter that silently bypassed the checkpoint's input normalization (Appendix~\ref{app:adapters}).
Table~\ref{tab:models} lists the models used in this paper.

\begin{table}[t]
\centering
\footnotesize
\setlength{\tabcolsep}{4pt}
\resizebox{\textwidth}{!}{%
\begin{tabular}{@{}lllcccrr@{}}
\toprule
Model & Checkpoint & Architecture & Token & Head & Tier & Params & FLOPs obs. \\
\midrule
\textcolor{tfm}{TimesFM 2.5} & \texttt{google/timesfm-2.5-200m} & decoder-only & 32 & deterministic & 3 & 231M & 42.5\% \\
\textcolor{ctwo}{Chronos-2} & \texttt{amazon/chronos-2} & encoder (+group attn.) & 16 & deterministic & 3 & 119M & 97.5\% \\
\textcolor{sun}{Sundial} & \texttt{thuml/sundial-base-128m} & decoder + flow head & 16 & sampled & 2 & 128M & 22.4\% \\
\textcolor{bolt}{Chronos-Bolt} & \texttt{amazon/chronos-bolt-small} & encoder + patch head & 16 & deterministic & 2 & 48M & 79.3\% \\
\midrule
Chronos-T5 & \texttt{amazon/chronos-t5-*} & encoder--decoder & 1 & sampled & 3 & 8M--709M & 14.4\%$^{a}$ \\
Timer & \texttt{thuml/timer-base-84m} & decoder-only & 96 & deterministic & 2$^{b}$ & 84M & -- \\
Lag-Llama & \texttt{Lag-Llama} & decoder over lags & lags & sampled & 2 & -- & -- \\
\bottomrule
\end{tabular}}
\caption{Models with adapters. The first four form the case study of Section~\ref{sec:casestudy}.
``Token'' is the number of time steps per input token; ``FLOPs obs.''\ is the share of a full
forecast's FLOPs spent in captured blocks, measured with PyTorch's FLOP counter
(\texttt{budget/model\_budget.json}). Sundial's share is low because its sampled flow-matching head runs
outside the captured pass. $^{a}$Chronos-T5-Base; only the encoder is captured, so the sampled decoder
is never observed (\fid{DE-05}). $^{b}$Via the zero-code \texttt{generic\_hf} adapter.}
\label{tab:models}
\end{table}

\subsection{Time alignment across architectures}
\label{sec:alignment}

Cross-model analysis needs a common index. \tool{} measures each token's span from its impulse
response (\texttt{extraction/span\_discovery.py}): a perturbation injected at one time step should move
the tokens whose spans contain it, and the resulting profile is gated on peak-to-pedestal contrast and
contiguity. Given spans $I_j=[s_j,e_j)$ for tokens $j=1,\dots,n$ and windows
$W_w=[wL,(w+1)L)$ of width $L$ (32 by default), the overlap-pooling matrix is
\begin{equation}
  P_{wj} \;=\; \frac{|I_j \cap W_w|}{\sum_{j'} |I_{j'} \cap W_w|},
  \qquad \tilde h_w \;=\; \sum_j P_{wj}\, h_j ,
  \label{eq:pooling}
\end{equation}
and every cross-model analysis runs on the resulting $[\text{series},\text{window},\text{dim}]$ tensors
(Figure~\ref{fig:alignment}). A window that receives no token raises an error rather than producing a
zero row. During extraction an \emph{alignment gate} injects an impulse into each window and checks,
layer by layer, that the matching pooled window moves most. Because coarse tokens cannot resolve fine
windows (a 96-step token on 32-step windows can reach at most one third of diagonal hits), the gate
normalizes hits by their attainable ceiling (\fid{MP-02}). The impulse amplitude is self-calibrated per
checkpoint and context length.

\begin{figure}[t]
\centering
\begin{tikzpicture}[x=0.36mm,y=6mm,font=\scriptsize]
  \foreach \row/\name/\w/\col in {3/TimesFM (32)/32/tfm, 2/{Chronos-2, Sundial, Bolt (16)}/16/ctwo, 1/Chronos-T5 (1)/4/gray, 0/Timer (96)/96/sun}{
    \node[anchor=east] at (-4,\row) {\name};
    \pgfmathtruncatemacro{\n}{288/\w-1}
    \foreach \k in {0,...,\n}{
      \draw[fill=\col!25, draw=\col!80!black, line width=0.3pt] (\k*\w,\row-0.32) rectangle ({(\k+1)*\w},\row+0.32);
    }
  }
  \foreach \k in {0,...,8}{
    \draw[thick, accent] (\k*32,-1.0) rectangle ({(\k+1)*32},-0.45);
    \node[font=\tiny] at (\k*32+16,-0.72) {$W_{\k}$};
  }
  \node[anchor=east] at (-4,-0.72) {common windows ($L{=}32$)};
  \draw[-{Stealth}, accent] (80,-0.3) -- (80,-0.02);
  \draw[decorate,decoration={brace,amplitude=3pt,mirror}] (0,-1.15) -- (288,-1.15) node[midway,below=3pt] {time (first 288 context steps shown)};
\end{tikzpicture}
\caption{Architecture-agnostic time alignment. Tokens of each model cover measured, contiguous spans
of different widths (Chronos-T5's one-step tokens are drawn in groups of four); Equation~\eqref{eq:pooling}
pools them by overlap onto common windows, so that all cross-model statistics compare the same stretch of
time.}
\label{fig:alignment}
\end{figure}

\subsection{The stage pipeline and its gates}
\label{sec:stages}

An analysis is a directed pipeline of 19 stages (Figure~\ref{fig:stages}), each writing typed artifacts
into the run directory. Stages skip themselves when their artifacts exist, and each stage declares the
configuration keys its artifacts depend on. These are fingerprinted, so a changed setting marks exactly
the dependent artifacts stale and a stale artifact is refused. Three gates decide what runs and write
their decisions to JSON, so the report can name them as skip reasons: the \emph{tier gate} (does every
model support this stage?), the \emph{run-shape gate} (comparison-only stages are dropped for a single
model) and the \emph{routing gate} (models that are not time-localized receive behavioral analysis
only). Extraction captures the residual stream into a chunked zarr store in float16 and runs the
alignment gate. All interventions go through a single primitive (\texttt{hooks.token\_patch}), which
raises if the model silently splits a batch.

\begin{figure}[t]
\centering
\begin{tikzpicture}[
  node distance=2.2mm and 2.2mm,
  st/.style={draw, rounded corners=1.5pt, font=\scriptsize\sffamily, minimum width=13mm, minimum height=5.5mm, fill=white},
  arr/.style={-{Stealth[length=1.4mm]}, gray!70!black},
  band/.style={rounded corners=3pt, inner sep=3pt},
  blab/.style={font=\scriptsize\bfseries, anchor=west}]
  \node[st] (corpus) {corpus};
  \node[st, right=of corpus] (extract) {extract};
  \node[st, right=of extract] (budget) {budget};
  \node[st, right=of budget] (frontend) {frontend};
  \node[st, right=of frontend] (screen) {layer\_screen};
  \node[st, below=5mm of corpus] (l0) {l0};
  \node[st, right=of l0] (internals) {internals};
  \node[st, right=of internals] (lens) {lens};
  \node[st, below=5mm of l0] (l1) {l1};
  \node[st, right=of l1] (l2) {l2};
  \node[st, right=of l2] (cluster) {cluster};
  \node[st, below=5mm of l1] (l3) {l3};
  \node[st, right=of l3] (attention) {attention};
  \node[st, right=of attention] (sae) {sae};
  \node[st, right=of sae] (concepts) {concepts};
  \node[st, below=5mm of l3] (exemplars) {exemplars};
  \node[st, right=of exemplars] (register) {register};
  \node[st, right=of register] (confirm) {confirm};
  \node[st, right=of confirm, fill=accent!15] (report) {report};
  \begin{scope}[on background layer]
    \node[band, fill=gray!10, fit=(corpus)(screen)] (b1) {};
    \node[band, fill=tfm!10, fit=(l0)(lens)] (b2) {};
    \node[band, fill=ctwo!10, fit=(l1)(cluster)] (b3) {};
    \node[band, fill=sun!12, fit=(l3)(concepts)] (b4) {};
    \node[band, fill=accent!10, fit=(exemplars)(report)] (b5) {};
  \end{scope}
  \node[blab] at ($(b1.east)+(2mm,0)$) {setup \& cost};
  \node[blab, align=left] at ($(lens.east)+(4mm,0)$) {behavioral / descriptive\\[-1pt]{\scriptsize\rmfamily\mdseries L0 accuracy, depth profile, lens}};
  \node[blab, align=left] at ($(cluster.east)+(4mm,0)$) {geometric / linear-translatable\\[-1pt]{\scriptsize\rmfamily\mdseries L1 CKA/RSA, L2 stitching, clusters}};
  \node[blab, align=left] at ($(b4.east)+(2mm,0)$) {causal within-model\\[-1pt]{\scriptsize\rmfamily\mdseries patching, heads, SAE ablation}};
  \node[blab, align=left] at ($(b5.east)+(2mm,0)$) {held-out confirmed};
  \draw[arr] (b1.south) -- (b2.north -| b1.south);
  \draw[arr] (b2.south -| b1.south) -- (b3.north -| b1.south);
  \draw[arr] (b3.south -| b1.south) -- (b4.north -| b1.south);
  \draw[arr] (b4.south -| b1.south) -- (b5.north -| b1.south);
  % gates
  \node[draw, dashed, font=\scriptsize, align=left, anchor=north west, text width=0.8\textwidth, fill=white] at ($(b5.south west)+(0,-3mm)$)
    {\textbf{Gates} (each writes JSON and a rendered reason): \emph{tier} (required adapter methods present?);
     \emph{run shape} (solo, pair or panel?); \emph{routing} (tokens time-localized?);
     \emph{fingerprints} (artifact computed under the current configuration?).};
\end{tikzpicture}
\caption{The 19-stage pipeline, grouped by the strongest evidence class each group can produce.
Stages below a gate that cannot run for a model are skipped for that model with the reason rendered in
the report. Only \texttt{confirm}, on the sealed private split, produces \ev{held-out confirmed}
evidence.}
\label{fig:stages}
\end{figure}

\subsection{The report: interpretability for non-specialists}
\label{sec:report}

The deliverable of a run is a single HTML file with no external dependencies
(\repo{examples/concept_atlas_v2/report.html}; Appendix~\ref{app:report} shows excerpts). It is designed
so that a reader with no background in interpretability can follow it, and so that an expert can audit
every number.

\begin{itemize}
  \item \textbf{Question-first structure.} After an ``at a glance'' block, seven parts pose plain
    questions (``Is the comparison fair?'', ``How well does each model forecast?'', ``Where does the
    forecast form?'', ``Do the models represent things the same way?'', ``What causally drives the
    forecast?'', ``Features and concepts'', ``Robustness and held-out confirmation''), and each stage
    section opens with a card stating its question, its method, what a good and a bad result look like,
    and what it cannot tell you.
  \item \textbf{Every figure explained.} Each figure has a caption plus a collapsible ``What does this
    mean?'' note giving its purpose, how to read it and its limitations; the example report contains
    361 such notes across its 19 sections. A glossary, a methods appendix and a gallery of known failure
    modes are generated from the code.
  \item \textbf{Derived, auditable verdicts.} The scorecard (33 rows in the example) is produced by a
    pure reduction over artifacts with no model names hard-coded: each row shows the measured value, the
    reference it is compared with (a null, a floor, a ceiling or a baseline) and the rule that turns the
    two into a verdict, so a reader can apply a different rule to the same numbers. A missing reference
    renders as ``not comparable''.
  \item \textbf{Typed findings.} Every claim is a \texttt{Finding} record with an identifier, stage,
    evidence class, technical text, a plain-language sentence and a caveat generated from the evidence
    class. The example run emits 126 findings (58 descriptive, 29 behavioral, 21 causal within-model,
    12 linear-translatable, 4 geometric, 2 illustrative), serialized to \texttt{report/findings.json}.
  \item \textbf{Nothing silently missing.} A stage that did not run renders as ``not run'' with its gate
    reason; a stage that raised renders as a visible ``section failed''. A multiplicity ledger lists every
    correction family (45 families and 3{,}599 comparisons in the example) and flags any family whose
    smallest attainable adjusted $p$-value cannot reach $\alpha$.
  \item \textbf{Reading levels.} A toggle switches between a headline view (fairness card, accuracy and
    confirmed findings only), the standard view and a methods view with every detail expanded.
\end{itemize}

\noindent The complete analysis is one command:
\begin{center}
\texttt{python run.py -{}-config configs/examples/large\_4model.yaml}
\end{center}
and smaller presets cover a CPU-only smoke test with mock models, a single real model and a pair
(\repo{tsfm_lens/configs/README.md}).
